\documentclass[10pt,a4paper]{article}
\usepackage{amsmath,amssymb,amsthm,graphicx}
\usepackage{amsmath,amssymb,amsthm,hyperref,tikz,enumerate,subfigure}
\newtheorem{theorem}{Theorem}[section]
\newtheorem{theo}[theorem]{Theorem}
\newtheorem{lemm}[theorem]{Lemma}
\newtheorem{prop}[theorem]{Proposition}
\newtheorem{defi}[theorem]{Definition}
\newtheorem{rema}[theorem]{Remark}
% Editorial retained-preamble fragment. The publisher class (cedram) defines the
% corollary environment used by the cached author source at lines 516 and 722;
% the layout adapter mirrors the other publisher theorem environments but not
% this one. This adds the matching environment with the same shared counter.
% No mathematical macro, symbol or assumption is introduced.
\newtheorem{coro}[theorem]{Corollary}
\DeclareMathSymbol{\leqslant}{\mathalpha}{AMSa}{"36}
\DeclareMathSymbol{\geqslant}{\mathalpha}{AMSa}{"3E}
\DeclareMathSymbol{\eset}{\mathalpha}{AMSb}{"3F}
\renewcommand{\leq}{\;\leqslant\;}
\renewcommand{\geq}{\;\geqslant\;}
%\newcommand{\dd}{\,\text{\rm d}}




\newcommand{\la}{\label}
\newcommand{\Om}{\Omega}


\newcommand{\e}{\varepsilon}
\newcommand{\h}{\frac{1}{2}}

\newcommand{\rhoN}{\rho^{\pi}}
\newcommand{\si}{\sigma}
\newcommand{\tc}{\, |\, }

\newcommand{\scalar}[2]{\langle #1, #2\rangle}



\DeclareMathOperator*{\union}{\bigcup}
\DeclareMathOperator*{\inter}{\bigcap}
\DeclareMathOperator{\var}{Var}
\newcommand{\tv}{\mathrm{TV}}
\newcommand{\ent}{\mathrm{Ent}}
\newcommand{\gap}{\mathrm{gap}}
\newcommand{\rBe}{\mathrm{Be}}
\newcommand{\IND}{\mathbf{1}}
\newcommand{\ind}{\mathbf{1}}

\newcommand{\cA}{\mathcal{A}}
\newcommand{\cB}{\mathcal{B}}
\newcommand{\cE}{\mathcal{E}}
\newcommand{\cL}{\mathcal{L}}
\newcommand{\cP}{\mathcal{P}}
\newcommand{\cS}{\mathcal{S}}
\newcommand{\cT}{\mathcal{T}}
\newcommand{\cU}{\mathcal{U}}
\newcommand{\cV}{\mathcal{V}}
\newcommand{\bbC}{\mathbb{C}}
\newcommand{\bbE}{\mathbb{E}}

\newcommand{\bbN}{\mathbb{N}}
\newcommand{\bbP}{\mathbb{P}}
\newcommand{\bbR}{\mathbb{R}}
\newcommand{\bbS}{\mathbb{S}}


%\DeclareMathOperator{\cov}{Cov}
%\DeclareMathOperator{\Hess}{Hess}
%\DeclareMathOperator{\dist}{dist}
%\DeclareMathOperator{\osc}{osc}
%\DeclareMathOperator{\Id}{Id}
%\DeclareMathOperator{\diag}{diag}
%\DeclareMathOperator{\supp}{supp}
%\DeclareMathOperator{\PI}{PI}
%\DeclareMathOperator{\Ent}{Ent}




\let\a=\alpha
\let\b=\beta
\let\d=\delta
\let\e=\varepsilon
\let\g=\gamma
\let\l=\lambda
\let\o=\omega
\let\r=\rho
\let\t=\tau
\let\z=\zeta
\let\G=\Gamma
\let\O=\Omega





%%%%%%%%%%%%%%%%%%%%%%%%%%%%%%%%%%%%%%%%%%%%%%%%%%%%%%%%%%%%%%%%%%%%%%%%%%%%%%%%%%%%%%%%%%%%%%%%%%%
\graphicspath{{./figures/}}

\newcommand*{\mk}{\mkern -1mu}
\newcommand*{\Mk}{\mkern -2mu}
\newcommand*{\mK}{\mkern 1mu}
\newcommand*{\MK}{\mkern 2mu}

\hypersetup{urlcolor=purple, linkcolor=blue, citecolor=red}

\newcommand*{\romanenumi}{\renewcommand*{\theenumi}{\roman{enumi}}}
\newcommand*{\Romanenumi}{\renewcommand*{\theenumi}{\Roman{enumi}}}
\newcommand*{\alphenumi}{\renewcommand*{\theenumi}{\alph{enumi}}}
\newcommand*{\Alphenumi}{\renewcommand*{\theenumi}{\Alph{enumi}}}
\let\oldtilde\tilde
\renewcommand*{\tilde}[1]{\mathchoice{\widetilde{#1}}{\widetilde{#1}}{\oldtilde{#1}}{\oldtilde{#1}}}
\let\oldhat\hat
\renewcommand*{\hat}[1]{\mathchoice{\widehat{#1}}{\widehat{#1}}{\oldhat{#1}}{\oldhat{#1}}}
\let\oldforall\forall
\renewcommand*{\forall}{\mathrel{\oldforall}}
\title{Kac chaos with explicit k/sqrt(N) and k/N total-variation rates for the canonical tensor product}
\author{Pietro Caputo, Daniel Parisi}
\date{doi:10.5802/ahl.219:published-2024; CC BY 4.0}
\begin{document}
\maketitle
\noindent\textit{Editorial note.} The statement is Theorem th:eqens of Section 2 (Local Central Limit Theorem and Chaos); only this theorem is counted here.\par\medskip
\noindent\textit{Editorial note.} The proof span ends at the closing of its own proof environment at line 717. The immediately following Corollary coro3 at 722-756, which bounds the expected total-variation distance of the empirical measure, is a separate consequence of this theorem and is deliberately not included; the remark at 719-720 explaining the f-g subtraction is also outside the span.\par\medskip
\noindent\textit{Editorial note.} The characteristic-function lemma 546-552 and the normalisation corollary coro2 516-530 are carried as context because the proof invokes them at 634 and 639; their own proofs are not part of this card.\par\medskip
\noindent\textit{Editorial note.} All displayed constants and rates are taken from the fixed publisher source unchanged.\par\medskip
\noindent\textit{Editorial note.} The cached author source uses the publisher's corollary environment at 516 (coro2) but the class file is not part of the deposit; the adapter mirrors the other theorem environments only. The one-line retained fragment sources/nonlinear-recombinations-ahl-2024-219/corollary-env.tex adds that same environment with the shared counter, and introduces no mathematical macro.\par\medskip
\noindent\textit{Allowed standard prerequisites (editorial).}\begin{itemize}
\item Kac chaos set-up and canonical tensor product. A symmetric mu\_N in P(O\textasciicircum{}N) is p-chaotic if every k-particle marginal P\_k mu\_N converges weakly to p\textasciicircum{}\{otimes k\}; gamma(p,rho\_N) is p\textasciicircum{}\{otimes N\} conditioned on Omega\_\{rho\_N\} and gamma\_N(p) = gamma(p,rho\textasciicircum{}pi\_N).. Fixes the objects compared in the two displays.
\item Multivariate CLT normalisation on the induced measure space. S\_N = sum\_\{j<=N\} xi(j) in \{0,...,N\}\textasciicircum{}K, with covariance V\_1, nondegenerate when det V\_1 != 0 and irreducible when every coordinate can be flipped with positive probability.. Supplies K, V\_1 and the admissible sequence convention; no rate is assumed.
\item Local CLT normalisation corollary. If mu is irreducible and the admissible sequence satisfies <rho\_N - pi, rho\_N - pi> = O(1/N), then mu\textasciicircum{}\{otimes N\}(Omega\_\{rho\_N\}) >= c N\textasciicircum{}\{-K/2\} for large N.. Gives the positive lower bound for the denominator of the Fourier quotient at line 634; the corollary itself is a consequence of the local CLT proved later in the paper.
\item Exponential decay of the characteristic function. For irreducible mu there is c = c(mu) > 0 with |psi(t)| <= exp(-c<t,t>) on [-pi,pi]\textasciicircum{}K.. Turns the far part of the Fourier integral into a Gaussian-type weight exp(-a<s,s>/N).
\end{itemize}\medskip
\begin{align}\label{rhoseq}
\rhoN_{i,x} :=\tfrac1N\,{\left\lfloor N \pi_{i,x}\right\rfloor}\,, \qquad i\in[n], \;x\in \{1,\,\dots,\,q_i\},
\end{align}
and we set $ \rhoN_{i,0}=1-\sum_{x=1}^{q_i}\rhoN_{i,x}$.

\subsection{Chaos and the propagation of chaos}

The chaos property is commonly defined as follows. A measure $\mu_N \in \cP(\O^{N})$ is \emph{symmetric} if it is invariant under any permutation of the $N$ particles. We write $ P_k\mu_N\in\cP(\O^k)$ for the corresponding $k$-particle marginal.


\begin{defi}[Kac's chaos or ``Boltzmann property'']\label{kacschaos}
A sequence $\mu_N \in \cP(\O^{N})$ of symmetric probabilities on $\Om^{N}$ is $\mu-$chaotic, for a given $\mu \in \cP(\O)$, if for any $k \in \bbN$ one has the weak convergence
\[
P_k\mu_N\longrightarrow \mu^{\otimes k}\,,\qquad N\to\infty.
\]
\end{defi}

A key step in implementing Kac's program is to construct a correspondence between probability measures on $\O$ and probability measures on $\O^N$. In our setting this can be formulated as follows.

\begin{defi}[Canonical tensor product]\label{def:canonical}
Given a probability measure $p \in \cP(\Om)$ and an admissible sequence of density vectors $\r_N$, we let
\[
\gamma(p,\r_N):=p^{\otimes N}\left(\,\cdot\,\middle|\,\Om_{\r_N}\right)
\]
be the tensor product of $p$ conditioned on $\Om_{\r_N}$. When $\r_N$ is given by $\rhoN$ as in~\eqref{rhoseq}, where the $\pi_{i,x}=p(\si_i=x)$ are the marginals of $p$, we use the notation $\g_N(p):= \gamma(p,\rhoN)$, and call it the \emph{canonical tensor product}.
\end{defi}

To avoid degeneracies we sometimes assume the following property.

\begin{align}\label{induce}
\si\in\O \mapsto \xi_{i,x} = \ind(\si_i=x)\,,\qquad i=1,\,\dots,\,n\,;\; x\in \{1,\,\dots,\,q_i\}.
\end{align}
That is, $\mu$ is the push forward of $p$ by the above map. Note that we did not include the indicator variable $ \ind(\si_i=0)$ since this is uniquely determined as the indicator of the event $\xi_{i,x}=0$ for all $x\in \{1,\,\dots,\,q_i\}$. When $q_i=1$ for all $i$, then $\O=\{0,1\}^n$ can be identified with $X$, $\si$ with $\xi$, and $\mu$ with $p$.


\subsection{Central limit theorem}

The next results are concerned with the behavior of the sum of independent copies $\xi(1),\,\dots,\,\xi(N)$ of a random variable $\xi$ with values in $X$ and distribution $\mu\in\cP(X)$:
\[
S_N = \sum_{j=1}^N\xi(j).
\]
Thus, $S_N$ is a random vector in $\{0,\,\dots,\,N\}^K$. We use the notation $\scalar{t}{s}=\sum_{i,x}t_{i,x}s_{i,x}$ if $t$ and $s$ are indexed by $i=1,\,\dots,\,n$ and $x=1,\,\dots,\,q_i$. We call $V_1$ the covariance matrix of $\mu$,
\[
V_1(i,x;i',x')=\mu(\xi_{i,x}\xi_{i',x'})-\mu(\xi_{i,x})\mu(\xi_{i',x'}),\qquad i=1,\dots,n\,,\; x=1,\,\dots,\,q_i
\]
Thus $V_1$ is a symmetric nonnegative definite $K\times K$ matrix. If $\det V_1\neq 0$ we say that $\mu$ is \emph{nondegenerate}. The central limit theorem asserts that if $\mu $ is nondegenerate, then as $N\to\infty$ one has the weak convergence
\begin{align}\label{eq:CLT}
\frac1{\sqrt N} \;V_1^{-1/2}\left(S_N - \mu^{\otimes N}(S_N)\right)\;
\longrightarrow\; N(0,\ind_K)\,,
\end{align}
where $\mu^{\otimes N}(S_N)\in[0,N]^K$ is the mean of the vector $S_N$ under the product measure $\mu^{\otimes N}$, and $\ind_K$ denotes the $K\times K$ identity matrix, so that $N(0,\ind_K)$ is the standard normal in $K$ dimensions. Note that when $\mu$ is induced by a measure $p\in\cP(\O)$ as described in~\eqref{induce}, then $\mu^{\otimes N}(S_N)_{i,x}=N\pi_{i,x}$ for all $i,x$, where $\pi_{i,x}$ are the marginals of $p$.

The statement~\eqref{eq:CLT} clearly requires that $\mu$ is nondegenerate. However, one can obtain similar statements in the case of degenerate measures, provided one reduces to the nondegenerate modes by eliminating the degenerate ones. More precisely, one can take the eigenvectors of $V_1$ with nonzero eigenvalues as the new variables. A simple example is obtained if e.g.\ $q_i\equiv 1$ and $\mu$ gives probability $1/2$ to all $1$'s and probability $1/2$ to all $0$'s. Here one simply removes all variables but one.

\subsection{Local central limit theorem}

We will need a local version of the central limit theorem. For this we assume the following stronger notion of nondegeneracy, which we refer to as irreducibility.

\begin{defi}\label{def:irr}
A measure $\mu\in\cP(X)$ is called \emph{irreducible} if for all $i=1,\,\dots, \,n$, for all $x\in\{0,\,\dots,\,q_i\}$, there exists $\xi\in X$ such that $\mu(\xi)$ and $\mu(\xi^{(i,x)})$ are both positive, where $\xi^{(i,x)}$ denotes the vector $\xi$ with the $(i,x)^{\rm th}$ coordinate flipped, that is $\xi^{(i,x)}_{j,y}=\xi_{j,y}$ for all $(j,y)\neq (i,x)$, and $\xi^{(i,x)}_{i,x}=1-\xi_{i,x}$.
\end{defi}

\begin{prop}\label{prop:LCLT}
Suppose $\mu\in\cP(X)$ is irreducible. Then there exists a finite constant $C=C(\mu)$ such that for all $N\in\bbN$,
\begin{align}\label{eq:LCLT}
\max_{M_N}\;
\left |\;\mu^{\otimes N}\left(S_N=M_N\right) \;-\; \frac{e^{-\frac12\,\scalar{z_N}{z_N}}}{(2\pi N) ^{K/2}
\sqrt {\det V_1} }\;
\right| \leq \frac{C}{ N^{(K+1)/2}},
\end{align}
where
\[
z_N:=\frac1{\sqrt N} \;V_1^{-1/2}\left(M_N - \mu^{\otimes N}(S_N)\right),
\]
and the maximum is over all possible values $M_N\in\{0,\,\dots,\,N\}^K$.
\end{prop}
Noting that $\O_{\r_N}=\{S_N=M_N\}$ with $M_N=N\r_N$, and that in this case
\[
\scalar{z_N}{z_N}
= N \left\langle\r_N-\pi,V_1^{-1}(\r_N-\pi)\right\rangle,
\]
the following is an immediate corollary of Proposition~\ref{prop:LCLT}.

\begin{coro}\label{coro2}
Suppose $\mu\in\cP(X)$ is irreducible, and let $\r_N$ be an admissible sequence such that
\begin{align}\label{hypo1}
\scalar{\r_N-\pi}{\r_N-\pi} =O(1/N),
\end{align}
where $\pi = \mu(\xi)$ is the vector of the expected values of $\mu$. Then there exists a constant $c=c(\mu)>0$ such that for $N$ sufficiently large
\begin{align}\label{lim1}
\mu^{\otimes N}\left(\Om_{\r_N}\right)\geq \frac{c}{N^{K/2}}.
\end{align}
In particular,
\begin{align}\label{loglim}
\lim_{N\,\to\,\infty} \frac1N\,\log
\mu^{\otimes N}\left(\Om_{\r_N}\right)= 0.
\end{align}
\end{coro}

\begin{lemm}\label{charfunc}
Suppose $\mu$ is irreducible. Then there exists a constant $c=c(\mu)>0$ such that the characteristic function $\psi(t)=\mu(e^{i\scalar{t}{\xi}})$, $t\in\bbR^K$, satisfies
\begin{align*}
|\psi(t)|\leq e^{-c
\scalar{t}{t} }\,,\qquad \text{for all\;\;} t\in[-\pi,\pi]^K.
\end{align*}
\end{lemm}

\subsection{Kac chaos}

Recall the definition of $\g(p,\r_N)$ and of the canonical tensor product $\g_N(p)$ in Definition~\ref{def:canonical}.

\begin{theo}\label{th:eqens}
Suppose $p\in\cP(\O)$ is irreducible and let $\r_N$ be an admissible sequence such that
\begin{align}\label{hypo10}
\scalar{\r_N-\pi}{\r_N-\pi} =O(1/N).
\end{align}
Then for all $k=1,\,\dots,\,N$,
\begin{align}\label{eq:chaos00}
\left\|P_k\g(p,\r_N)-p^{\otimes k}\right\|_\tv \leq \frac{C\,k}{\sqrt N},
\end{align}
for some constant $C=C(p)$. Moreover, when $\r_N=\rhoN$, the canonical tensor product $\g_N(p)$ satisfies the stronger estimate
\begin{align}\label{eq:chaos000}
\left\|P_k\g_N(p)-p^{\otimes k}\right\|_\tv \leq \frac{C\,k}{ N}.
\end{align}
\end{theo}

\begin{proof}
We prove~\eqref{eq:chaos00} first, and then show how to obtain~\eqref{eq:chaos000}. Let $\hat \xi_{i,x}(j)=\xi_{i,x}(j)-(\r_N)_{i,x}$. We use the shorthand notation $\g_N=\g(p,\r_N)$ and $\mu_N=p^{\otimes N}$. For any $f:\O^N\mapsto\bbR$ we have
\begin{align}\label{eq:chao1}
\g_N(f)-\mu_N(f)=\frac{\mu_N(f(\ind_{\O_{\r_N}}-\mu_N(\O_{\r_N}))}{\mu_N(\O_{\r_N})},
\end{align}
Since $\O_{\r_N}=\{S_N=N\r_N\}$, using the Fourier transform we write
\begin{align}\label{eq:chao2}
\mu_N(\O_{\r_N})=\frac1{(2\pi)^K}\int_{[-\pi,\pi]^K}d t \,\mu_N\left(e^{i\left\langle t,\hat S_N\right\rangle}\right),
\end{align}
where $\hat S_N=\sum_{j=1}^N\hat\xi(j)$. Set $V_N:=N V_1$. The product structure of $\mu_N$ and the change of variables $s= V_N^{1/2} t=\sqrt N V_1^{1/2} t\,$ imply
\begin{align}\label{eq:chao3}
\mu_N(\O_{\r_N})=\frac1{B_N(2\pi)^K}\int_{Q_{N,K}}d s \,\mu_N\left(e^{i\left\langle V_N^{-1/2} s,\hat\xi(1)\right\rangle}\right)^N,
\end{align}
where $Q_{N,K}=V_N^{1/2}[-\pi,\pi]^K$ and $B_N = \sqrt {\det V_N}=N^{K/2}\sqrt {\det V_1}$.

In the same way, for any $f=f(\xi(1),\,\dots,\,\xi(k))$, we have
\begin{multline}\label{eq:chao4}
\mu_N(f\,\ind_{\O_{\r_N}})\\
=\frac1{B_N(2\pi)^K}\int_{Q_{N,K}}d s \,\mu_N\left(e^{i\left\langle V_N^{-1/2} s,\hat\xi(1)\right\rangle}\right)^{N-k}\mu_N\left(f \,e^{i\left\langle V_N^{-1/2} s,\hat S_k\right\rangle}\right).
\end{multline}
In conclusion, we have
\begin{align}
\g_N(f)-\mu_N(f) &=\frac{\int_{Q_{N,K}}d s \,\psi_N(s)^{N-k}\mu_N\left(f; \,e^{i\left\langle V_N^{-1/2} s,\hat S_k\right\rangle}\right)}{
\int_{Q_{N,K}}d s \,\psi_N(s)^{N}\,},\label{eq:chao5}
\\
\intertext{where}
\psi_N(s) &=\mu_N\left(e^{i\left\langle V_N^{-1/2} s,\hat\xi(1)\right\rangle}\right),\nonumber
\end{align}
and we use the notation $\mu_N(f; g)=\mu_N(fg)-\mu_N(f)\mu_N(g)$ for the covariance of $f,g$. From Corollary~\ref{coro2} we known that~\eqref{eq:chao3} is at least $cN^{-K/2}$, and thus the denominator in~\eqref{eq:chao5} is at least some constant $c'>0$.
Therefore, it suffices to show that the numerator is bounded by
\begin{align}\label{eq:chao6}
\int_{Q_{N,K}}d s \,|\psi_N(s)|^{N-k}\left|\mu_N\left(f ;\,e^{i\left\langle V_N^{-1/2} s,\hat S_k\right\rangle}\right)\right|\leq C \,|f|_\infty\,\frac{k}{\sqrt N}.
\end{align}
From Lemma~\ref{charfunc} we know that $|\psi_N(s)|\leq e^{-a \scalar{s}{s}/N}$ for some constant $a=a(p)>0$. Notice that we can assume without loss of generality that $k\leq N/2$, since otherwise the result~\eqref{eq:chaos00} is trivial. Thus $|\psi_N(s)|^{N-k} \leq e^{-a \scalar{s}{s}/2}$ and it is sufficient to show that
\begin{align}\label{eq:chao7}
\left|\mu_N\left(f; \,e^{i\left\langle V_N^{-1/2} s,\hat S_k\right\rangle}\right)\right|\leq C \,|f|_\infty\,\frac{k}{\sqrt N}\,\scalar{s}{s},
\end{align}
for all $s\in Q_{N,K}$. Recalling that $|e^{i\theta}-1|\leq |\theta|$, $\theta\in\bbR$, and using Schwarz' inequality,
\begin{align}\label{eq:chao71}
\left|\mu_N\left(f; \,e^{i\left\langle V_N^{-1/2} s,\hat S_k\right\rangle}
\right)\right|
\leq \frac1{\sqrt N} \,|f|_\infty\,
\mu_N\left(\left\langle V_1^{-1/2} s,\hat S_k\right\rangle^2\right).
\end{align}
Now we observe that
\begin{multline}\label{eq:chao72}
\mu_N\left(\left\langle V_1^{-1/2} s,\hat S_k\right\rangle^2\right) = k^2\scalar{V_1^{-1/2} s}{\pi-\r_N}^2 + k \scalar{s}{s}
\\
 \leq k^2\left\langle s,V_1^{-1}s\right\rangle \left\langle\pi-\r_N,\pi-\r_N\right\rangle + k \scalar{s}{s} \leq Ck \scalar{s}{s},
\end{multline}
where we use $\mu_N(\hat S_k)=k\,(\pi-\r_N)$, the independence of the $\hat\xi(j)$, and $\scalar{\pi-\r_N}{\pi-\r_N}\leq C/k$ which follows from the assumption $\scalar{\pi-\r_N}{\pi-\r_N}=O(1/N)$. This proves~\eqref{eq:chaos00}.

To prove~\eqref{eq:chaos000}, note that it is sufficient to prove~\eqref{eq:chao7} with $\sqrt N$ replaced by $N$ in the right hand side. For this, we are going to use the fact that $\scalar{\pi-\r_N}{\pi-\r_N}=O(1/N^2)$ when $\r_N=\rhoN$, see~\eqref{rhoseq}. Let us first consider the function $\tilde f = f-g$, where
\begin{align}\label{eq:chao8}
g(\xi(1),\,\dots,\,\xi(k))=\frac1k\left\langle V_1^{-1}v,\hat S_k\right\rangle,\qquad v = \mu_N\left(f;\hat S_k\right).
\end{align}
The function $g$ can be seen as a linear approximation of $f$. Notice that
\begin{align}\label{eq:chao9}
\mu_N\left(f-g\, \,;\,\left\langle V_N^{-1/2} s,\hat S_k\right\rangle\right)=0.
\end{align}
Indeed, by independence $\mu_N((\hat S_k)_{i,x};(\hat S_k)_{j,y}) = k V_1(i,x;j,y) $, and therefore for all $s$,
\begin{align}\label{eq:chao10}
\mu_N\left(g \,;\,\left\langle V_N^{-1/2} s,\hat S_k\right\rangle\right) = \left\langle v,V_N^{-1/2} s\right\rangle =\mu_N\left(f\,;\,\left\langle V_N^{-1/2} s,\hat S_k\right\rangle\right).
\end{align}
Recalling that $|e^{i\theta}-1-i\theta|=|R(\theta)|\leq \frac12\theta^2$, $\theta\in\bbR$, from~\eqref{eq:chao9} we have
\begin{multline}\label{eq:chao71x}
\left|\mu_N\left(f-g\,; \,e^{i\left\langle V_N^{-1/2} s,\hat S_k\right\rangle}
\right)\right|^2\\
\leq \frac1{N^2}\,\left(\var(f) + \var(g)\right)\mu_N\left(\left\langle V_1^{-1/2} s,\hat S_k\right\rangle^4\right)\,,
\end{multline}
where we use the inequality $|\mu_N(f-g;R)|^2\leq 2(\var(f) + \var(g))\mu_N(R^2)$, and we use $\var$ for the variance w.r.t.\ $\mu_N$. Next, observe that $\var(f)\leq |f|_\infty^2$ and
\[
\var(g) = \frac1k \left\langle V^{-1}v,v\right\rangle \leq C \var(f) \leq C |f|_\infty^2.
\]
We are going to show that
\begin{align}\label{eq:chao7a7}
\mu_N\left(\left\langle V_1^{-1/2} s,\hat S_k\right\rangle^4\right)
\leq C k^2\scalar{s}{s}^2\,.
\end{align}
Suppose for a moment that~\eqref{eq:chao7a7} holds. Then we conclude that
\begin{align}\label{eq:chao7a}
\left|\mu_N\left(f-g\,; \,e^{i\left\langle V_N^{-1/2} s,\hat S_k\right\rangle}\right)\right|\leq C \,|f|_\infty\,\frac{k}{ N}\,\scalar{s}{s}.
\end{align}
This proves $|\g_N(\tilde f)-\mu_N(\tilde f)|\leq C \,|f|_\infty\,k/N$. However, noting that $\g_N(g)=0$, we have
\begin{align}\label{eq:chaoa1}
|\g_N(f)-\mu_N(f)|\leq \left|\g_N(\tilde f)-\mu_N(\tilde f)\right|+|\mu_N(g)|.
\end{align}
The desired conclusion $|\g_N(f)-\mu_N(f)|\leq C \,|f|_\infty\,k/N$ then follows from
\begin{align}\label{eq:chaoap1}
\begin{split}
|\mu_N(g)|&\leq \left|\left\langle V_1^{-1}v,\pi-\rhoN\right\rangle\right|\\
&\leq \left\langle V_1^{-1}v,V_1^{-1}v\right\rangle^{1/2}\scalar{\pi-\rhoN}{\pi-\rhoN}^{1/2} \leq \frac{C|f|_\infty k}{N},
\end{split}
\end{align}
where we use $\scalar{\pi-\rhoN}{\pi-\rhoN}=O(1/N^2)$, and $\scalar{V_1^{-1}v}{V_1^{-1}v}\leq Ck|f|_\infty$.

Thus, it remains to prove~\eqref{eq:chao7a7}. Notice that it is sufficient to prove
\begin{align}\label{eq:chao7a8}
\mu_N\left(\left\langle\hat S_k,\hat S_k\right\rangle^2\right)
\leq C k^2\,.
\end{align}
We write $\tilde \xi = \hat \xi - \mu(\hat \xi)$ and $\tilde S_k=\hat S_k - \mu_N(\hat S_k)$ for the corresponding sums. Then one checks that
\begin{multline}\label{eq:chao7a9}
\left\langle\hat S_k,\hat S_k\right\rangle^2
\leq \left\langle\tilde S_k,\tilde S_k\right\rangle^2 \\
+ Ck\left\langle\tilde S_k,\tilde S_k\right\rangle^{3/2}\left\langle\pi-\rhoN,\pi-\rhoN\right\rangle^{1/2}
+ Ck^2\left\langle\tilde S_k,\tilde S_k\right\rangle
\left\langle\pi-\rhoN,\pi-\rhoN\right\rangle \\
+ Ck^3\left\langle\tilde S_k,\tilde S_k\right\rangle^{1/2}\left\langle\pi-\rhoN,\pi-\rhoN\right\rangle^{3/2} + Ck^4\left\langle\pi-\rhoN,\pi-\rhoN\right\rangle^2\,.
\end{multline}
Since $\scalar{\pi-\rhoN}{\pi-\rhoN}=O(1/N^2)$ we can restrict to prove~\eqref{eq:chao7a8} for $\tilde S_k$ instead of $\hat S_k$. Since $\tilde \xi$ are centered the estimate follows easily by expanding $\scalar{\tilde S_k}{\tilde S_k}^2$ and observing that the dominant terms are of the form $\scalar{\tilde\xi(i)}{\tilde\xi(j)}^2$, and their contribution is of order $k^2$.
\end{proof}

\begin{thebibliography}{99}
\bibitem{Bertini_etal} Basile, G.; Benedetto, D.; Caglioti, E.; Bertini, L. Large deviations for a binary collision model: energy evaporation, Math. Eng., Volume 5 (2023) no. 1, 1, p. 12
\bibitem{barthe_et_al} Barthe, F.; Cordero-Erausquin, D.; Ledoux, M.; Maurey, B. Correlation and Brascamp–Lieb inequalities for Markov semigroups, Int. Math. Res. Not., Volume 2011 (2011) no. 10, pp. 2177-2216
\bibitem{bonetto2017entropy} Bonetto, F.; Geisinger, A.; Loss, M.; Ried, T. Entropy decay for the Kac evolution, Commun. Math. Phys., Volume 363 (2018) no. 3, pp. 847-875
\bibitem{BT} Bobkov, S. G.; Tetali, P. Modified logarithmic Sobolev inequalities in discrete settings, J. Theor. Probab., Volume 19 (2006) no. 2, pp. 289-336
\bibitem{asymm} Caputo, P.; Martinelli, F. Asymmetric diffusion and the energy gap above the 111 ground state of the quantum XXZ model, Commun. Math. Phys., Volume 226 (2002) no. 2, pp. 323-375
\bibitem{entprod} Caputo, P.; Sinclair, A. Entropy production in nonlinear recombination models, Bernoulli, Volume 24 (2018) no. 4B, pp. 3246-3282
\bibitem{carlenwild} Carlen, E. A.; Carvalho, M. C.; Gabetta, E. Central limit theorem for Maxwellian molecules and truncation of the Wild expansion, Commun. Pure Appl. Math., Volume 53 (2000) no. 3, pp. 370-397
\bibitem{entfound} Carlen, E. A.; Carvalho, M. C.; Le Roux, J.; Loss, M.; Villani, C. Entropy and chaos in the Kac model, Kinet. Relat. Models, Volume 3 (2010) no. 1, pp. 85-122
\bibitem{revchaos} Chaintron, L.-P.; Diez, A. Propagation of chaos: a review of models, methods and applications. I: Models and methods, Kinet. Relat. Models, Volume 15 (2022) no. 6, pp. 895-1015
\bibitem{cattiaux} Cattiaux, P.; Guillin, A.; Malrieu, F. Probabilistic approach for granular media equations in the non-uniformly convex case, Probab. Theory Relat. Fields, Volume 140 (2008) no. 1-2, pp. 19-40
\bibitem{Chung} Chung, K. L. A course in probability theory, Academic Press Inc., 2000
\bibitem{Virag} Csóka, E.; Harangi, V.; Virág, B. Entropy and expansion, Ann. Inst. Henri Poincaré, Probab. Stat., Volume 56 (2020) no. 4, pp. 2428-2444
\bibitem{elementary} Durmus, A.; Eberle, A.; Guillin, A.; Zimmer, R. An elementary approach to uniform in time propagation of chaos, Proc. Am. Math. Soc., Volume 148 (2020) no. 12, pp. 5387-5398
\bibitem{dolera_etal} Dolera, E.; Gabetta, E.; Regazzini, E. Reaching the best possible rate of convergence to equilibrium for solutions of Kac’s equation via central limit theorem, Ann. Appl. Probab., Volume 19 (2009) no. 1, pp. 186-209
\bibitem{DelMoralGuionnet} Del Moral, P.; Guionnet, A. On the stability of interacting processes with applications to filtering and genetic algorithms, Ann. Inst. Henri Poincaré, Probab. Stat., Volume 37 (2001) no. 2, pp. 155-194
\bibitem{miclo} Del Moral, P.; Miclo, L. Genealogies and increasing propagation of chaos for Feynman–Kac and genetic models, Ann. Appl. Probab., Volume 11 (2001) no. 4, pp. 1166-1198
\bibitem{DiacShash} Diaconis, P.; Shahshahani, M. Generating a random permutation with random transpositions, Z. Wahrscheinlichkeitstheor. Verw. Geb., Volume 57 (1981), pp. 159-179
\bibitem{DS} Diaconis, P.; Saloff-Coste, L. Logarithmic Sobolev inequalities for finite Markov chains, Ann. Appl. Probab., Volume 6 (1996) no. 3, pp. 695-750
\bibitem{erbar} Erbar, M.; Fathi, M.; Schlichting, A. Entropic curvature and convergence to equilibrium for mean-field dynamics on discrete spaces, ALEA, Lat. Am. J. Probab. Math. Stat., Volume 17 (2020) no. 1, pp. 445-471
\bibitem{einav2011villani} Einav, A. On Villani’s conjecture concerning entropy production for the Kac master equation, Kinet. Relat. Models, Volume 4 (2011) no. 2, pp. 479-497
\bibitem{Gei} Geiringer, H. On the probability theory of linkage in Mendelian heredity, Ann. Math. Stat., Volume 15 (1944), pp. 25-57
\bibitem{Goel} Goel, S. Modified logarithmic Sobolev inequalities for some models of random walk, Stochastic Processes Appl., Volume 114 (2004) no. 1, pp. 51-79
\bibitem{GQ} Gao, F.; Quastel, J. Exponential decay of entropy in the random transposition and Bernoulli–Laplace models, Ann. Appl. Probab., Volume 13 (2003) no. 4, pp. 1591-1600
\bibitem{grunbaumchaos} Grünbaum, F. A. Propagation of chaos for the Boltzmann equation, Arch. Ration. Mech. Anal., Volume 42 (1971), pp. 323-345
\bibitem{onkacschaos} Hauray, M.; Mischler, S. On Kac’s chaos and related problems, J. Funct. Anal., Volume 266 (2014) no. 10, pp. 6055-6157
\bibitem{KAC} Kac, M. Foundations of kinetic theory, Proceedings of the Third Berkeley Symposium on Mathematical Statistics and Probability, 1954–1955, vol. III, University of California Press (1956), pp. 171-197
\bibitem{lacker} Lacker, D. Hierarchies, entropy, and quantitative propagation of chaos for mean field diffusions, Probab. Math. Phys., Volume 4 (2023) no. 2, pp. 377-432
\bibitem{Lyubich} Lyubich, Y. I. Mathematical structures in population genetics, Biomathematics (Berlin), 22, Springer, 1992 (translated from the Russian by D. Vulis and A. Karpov. Edited by E. Akin)
\bibitem{Martinez} Martínez, S. A probabilistic analysis of a discrete-time evolution in recombination, Adv. Appl. Math., Volume 91 (2017), pp. 115-136
\bibitem{mckean1966speed} McKean, H. P. jun. Speed of approach to equilibrium for Kac’s caricature of a Maxwellian gas, Arch. Ration. Mech. Anal., Volume 21 (1966), pp. 343-367
\bibitem{mckean} McKean, H. P. jun. An exponential formula for solving Boltzmann’s equation for a Maxwellian gas, Comb. Theory, Volume 2 (1967), pp. 358-382
\bibitem{kacsprog} Mischler, S.; Mouhot, C. Kac’s program in kinetic theory, Invent. Math., Volume 193 (2013) no. 1, pp. 1-147
\bibitem{anewapp} Mischler, S.; Mouhot, C.; Wennberg, B. A new approach to quantitative propagation of chaos for drift, diffusion and jump processes, Probab. Theory Relat. Fields, Volume 161 (2015) no. 1-2, pp. 1-59
\bibitem{Nagylaki} Nagylaki, T. Introduction to theoretical population genetics, Biomathematics (Berlin), 21, Springer, 2013
\bibitem{Petrov} Petrov, V. V. Sums of independent random variables, Walter de Gruyter, 2022
\bibitem{review_discrete} Platkowski, T.; Illner, R. Discrete velocity models of the Boltzmann equation: a survey on the mathematical aspects of the theory, SIAM Rev., Volume 30 (1988) no. 2, pp. 213-255
\bibitem{Reza} Rezakhanlou, F. Propagation of chaos for particle systems associated with discrete Boltzmann equatio, Stochastic Processes Appl., Volume 64 (1996) no. 1, pp. 55-72
\bibitem{Sinetal2} Rabani, Y.; Rabinovich, Y.; Sinclair, A. A computational view of population genetics, Random Struct. Algorithms, Volume 12 (1998) no. 4, pp. 313-334
\bibitem{Sinetal} Rabinovich, Y.; Sinclair, A.; Wigderson, A. Quadratic dynamical systems, 33rd annual symposium on Foundations of computer science (FOCS)Proceedings, IEEE (1992), pp. 304-313
\bibitem{Baakeetal} Salamat, M.; Baake, M.; Baake, E. The general recombination equation in continuous time and its solution, Discrete Contin. Dyn. Syst., Volume 36 (2016) no. 1, pp. 63-95
\bibitem{sznitman} Sznitman, A.-S. Topics in propagation of chaos, Calcul des probabilités, Ec. d’Été, Saint-Flour/Fr (Lecture Notes in Mathematics), Volume 1464, Springer (1989), pp. 165-251
\bibitem{tanaka} Tanaka, H. Probabilistic treatment of the Boltzmann equation of Maxwellian molecules, Z. Wahrscheinlichkeitstheor. Verw. Geb., Volume 46 (1978), pp. 67-105
\bibitem{Vil} Villani, C. Cercignani’s conjecture is sometimes true and always almost true, Commun. Math. Phys., Volume 234 (2003) no. 3, pp. 455-490
\bibitem{Wild} Wild, E. On Boltzmann’s equation in the kinetic theory of gases, Proc. Camb. Philos. Soc., Volume 47 (1951) no. 3, pp. 602-609
\end{thebibliography}
\end{document}
